\section{PDF 的最终结果}\label{sec:final_results}

\begin{figure}[tbp]
\includegraphics[width=.48\textwidth]{images/unpolarized_different_momentum}
\caption{\label{fig:matchingPDFallPz} 匹配后非极化同位旋矢量 PDF 对核子 boost 动量的依赖：绿色点线与蓝色实线分别对应核子动量 $P_z$ 为 1.8 与 2.3 GeV。傅里叶变换的截断选为 $zP^z\sim 12$（$P_z=1.8$ GeV 时 $z=15a$，$P_z=2.3$ GeV 时 $z=12a$）。}
\end{figure}

采用导数法 FT 以及 $p_z^R=0$ GeV、$\mu_R=3.2$ GeV 的匹配，我们在图~\ref{fig:matchingPDFallPz} 中展示了带统计不确定度的结果对核子 boost 动量的依赖。正如从图~\ref{fig:renormalized_h-tilde} 中准 PDF 矩阵元素的一致性所能预期的，二者彼此一致。

最后，我们在图~\ref{fig:finalPDF} 中给出 PDF 结果及其与全局分析~\cite{Dulat:2015mca,Ball:2017nwa,Harland-Lang:2014zoa}的比较。从图中可见，我们的结果在大 $x$ 区符合得较好，但在小 $x$ 区存在显著差异，主要来自 FT 截断方法以及 $p_z^R$ 依赖的系统不确定度。

\begin{figure}[tbp]
\includegraphics[width=.48\textwidth]{images/unpolarized_Pz=5_with_pheno}
\caption{\label{fig:finalPDF} 由核子动量 $P_z=2.3$~GeV 的 RI/MOM 准 PDF 计算的 $\mu=2$~GeV PDF 结果：与 CT14nnlo（90CL）~\cite{Dulat:2015mca}、NNPDF3.1（68CL）~\cite{Ball:2017nwa} 及 MMHT2014（68CL）~\cite{Harland-Lang:2014zoa} 比较。我们的结果在不确定度范围内与全局分析相符。}
\end{figure}

